% eksperymentalne tłumaczenie na włoski

%

\index{wzór!Herona} % lang-pl

\subsection{Pole trójkąta, wzór Herona} % lang-pl
\index{formula!di Erone} % lang-it
\subsection{Area del triangolo, formula di Erone} % lang-it


W tej podsekcji znajdziemy różne wzory na pole trójkąta. % lang-pl
In questa sottosezione presenteremo diverse formule per l'area di un triangolo. % lang-it

\begin{proposition}
    Pole trójkąta o podstawie długości $a$ oraz wysokości opuszczonej na tę podstawę długości $h$ wyraża się wzorem % lang-pl
    
    L'area di un triangolo con base di lunghezza $a$ e altezza relativa a tale base di lunghezza $h$ è data dalla formula % lang-it
    \begin{equation}
        S_\triangle = \frac 1 2 a h.
    \end{equation}
    Jeśli kąt naprzeciw podstawy ma miarę $\alpha$, a ramiona długości $b, c$, to zachodzi też % lang-pl
    
    Se l'angolo opposto alla base misura $\alpha$ e gli altri due lati hanno lunghezze $b$ e $c$, allora vale anche % lang-it
    \begin{equation}
        S_\triangle = \frac 1 2 bc \sin \alpha.
    \end{equation}
\end{proposition}

Pisze o tym Guzicki \cite[s. 255]{guzicki_2021}. % lang-pl

Se ne occupa Guzicki \cite[p.~255]{guzicki_2021}. % lang-it

\begin{proposition}[wzór Herona, 60 r.n.e.?] % lang-pl
\begin{proposition}[formula di Erone, 60 a.C.?] % lang-it
\label{prp_heron}%
    Pole trójkąta o bokach długości $a, b, c$ wyraża się wzorem % lang-pl
    
    L'area di un triangolo con lati di lunghezza $a$, $b$, $c$ è data dalla formula % lang-it
    \begin{equation}
        S = \sqrt{p(p-a)(p-b)(p-c)},
    \end{equation}
    gdzie $p = \frac 1 2 (a + b + c)$ oznacza połowę obwodu. % lang-pl
    
    dove $p = \frac 1 2 (a + b + c)$ denota il semiperimetro. % lang-it
\end{proposition}

Wynik przypiszemy Heronowi z Aleksandrii, poda dość skomplikowany dowód w~dziele \emph{,,Metrica''}. % lang-pl
Attribuiremo questo risultato a Erone di Alessandria, che ne fornirà una dimostrazione piuttosto complessa nell'opera \emph{Metrica}. % lang-it
\index[persons]{Erone di Alessandria}% % lang-it


\index[persons]{Heron z Aleksandrii}% % lang-pl
Coxeter \cite[s. 12]{coxeter_1991} (\cite[s. 28]{coxeter_1967}) przeczyta van der Waerdena \cite[s. 228, 277]{waerden_1961} i przez to stwierdzi, że wzór był już znany Archimedesowi dwieście lat wcześniej. % lang-pl
Coxeter \cite[p.~12]{coxeter_1991} (\cite[p.~28]{coxeter_1967}) leggerà van der Waerden \cite[pp.~228, 277]{waerden_1961} e concluderà che la formula era già nota ad Archimede duecento anni prima. % lang-it

% czemu jest cytowane angielskie i polskie wydanie?

\index[persons]{Archimedes}% % lang-pl
\index[persons]{Archimede}% % lang-it

Pompe w $\Delta_{93}^9$ pokaże, jak dojść do dowodu Sidneya Kunga ,,bez rachowania''. % lang-pl

Pompe mostrerà in $\Delta_{93}^9$ come ottenere la dimostrazione di Sidney Kung ``senza fare calcoli''. % lang-it
% https://www.deltami.edu.pl/1993/09/jak-udowodnic-wzor-herona/

\index[persons]{Kung, Sidney}%

Guzicki \cite[s. 165-169]{guzicki_2021} wyprowadzi go z twierdzenia Pitagorasa i~wspomni, że prosty geometryczny dowód poda(ł) też Euler. % lang-pl
Guzicki \cite[pp.~165--169]{guzicki_2021} lo dedurrà dal teorema di Pitagora e ricorderà che una semplice dimostrazione geometrica fu data anche da Eulero. % lang-it
\index[persons]{Euler, Leonhard}% % lang-it


\index[persons]{Euler, Leonhard}% % lang-pl
\index{twierdzenie!Pitagorasa} % lang-pl
\index{teorema!di Pitagora} % lang-it


Wreszcie Bogdańska, Neugebauer \cite[s. 92]{neugebauer_2018} skorzystają z twierdzenia kosinusów. % lang-pl
Infine Bogdańska e Neugebauer \cite[p.~92]{neugebauer_2018} utilizzeranno il teorema del coseno. % lang-it


Wzór Herona pojawi się wielokrotnie w innych numerach Delty: $\Delta_{14}^3$, $\Delta_{25}^{12}$ (jako pierwsz wpis na liście Marcusa Bakera stu dziesięciu różnych przepisów na pole trójkąta). % lang-pl
La formula di Erone comparirà più volte in altri numeri di Delta: $\Delta_{14}^3$, $\Delta_{25}^{12}$ (come prima voce nell'elenco di Marcus Baker di centodieci diverse formule per l'area del triangolo). % lang-it

% https://www.deltami.edu.pl/2014/03/heron-uogolniony/
% https://www.deltami.edu.pl/2025/12/heron-heron-po-trzykroc-heron/
% https://archive.org/details/jstor-1967384/page/n5/mode/2up
% https://archive.org/details/jstor-1967393

Istnieje częściowe uogólnienie wzoru Herona dla czworokątów: % lang-pl

Esiste una generalizzazione parziale della formula di Erone ai quadrilateri: % lang-it

\begin{proposition}[wzór Brahmagupty] % lang-pl
\begin{proposition}[formula di Brahmagupta] % lang-it
\label{brahmagupta_formula}%
\index{czworokąt!cykliczny}% % lang-pl
\index{wzór!Brahmagupty}% % lang-pl
    Pole czworokąta wypukłego, na którym można opisać okrąg, o bokach długości $a, b, c, d$ wyraża się wzorem % lang-pl
\index{quadrilatero!ciclico}% % lang-it
\index{formula!di Brahmagupta}% % lang-it
    
    L'area di un quadrilatero convesso ciclico con lati di lunghezza $a$, $b$, $c$, $d$ è data dalla formula % lang-it
    \begin{equation}
        S = \sqrt{(p-a)(p-b)(p-c)(p-d)},
    \end{equation}
    gdzie $p$ oznacza połowę obwodu. % lang-pl
    
    dove $p$ denota il semiperimetro. % lang-it
\end{proposition}

Dowód polega na zastosowaniu wzoru Herona dwa razy do trójkątów podobnych lub ponownym skorzystaniu z dobrodziejstw trygonometrii. % lang-pl
La dimostrazione consiste nell'applicare due volte la formula di Erone a opportuni triangoli oppure nel ricorrere nuovamente alla trigonometria. % lang-it


Założenie o wypukłości i cykliczności można pominąć, jak odkryją równocześnie Carl Bretschneider i Karl von Staudt: % lang-pl
Le ipotesi di convessità e ciclicità possono essere eliminate, come scopriranno indipendentemente Carl Bretschneider e Karl von Staudt: % lang-it

\index[persons]{Staudt@von Staudt, Karl}%
\index[persons]{Bretschneider, Carl}%

\begin{proposition}[wzór Bretschneidera, 1842] % lang-pl
\index{wzór!Bretschneidera}% % lang-pl
    Pole czworokąta o bokach długości $a, b, c, d$ i sumie miar dwóch przeciwległych kątów wewnętrznych $2 \varphi$ wyraża się wzorem % lang-pl
\begin{proposition}[formula di Bretschneider, 1842] % lang-it
\index{formula!di Bretschneider}% % lang-it
    
    L'area di un quadrilatero con lati di lunghezza $a$, $b$, $c$, $d$ e somma di due angoli opposti pari a $2\varphi$ è data dalla formula % lang-it
    \begin{equation}
        S = \sqrt{(p-a)(p-b)(p-c)(p-d) - abcd \cos^2  \varphi},
    \end{equation}
    gdzie $p$ oznacza połowę obwodu. % lang-pl
    
    dove $p$ denota il semiperimetro. % lang-it
\end{proposition}

Inny, równie przydatny wzór znajdzie Julian Coolidge \cite{coolidge_1939}: % lang-pl

Julian Coolidge \cite{coolidge_1939} troverà un'altra formula altrettanto utile: % lang-it

\index[persons]{Coolidge, Julian}%

pole czworokąta o bokach długości $a, b, c, d$ i przekątnych długości $e, f$ wyraża się wzorem % lang-pl

l'area di un quadrilatero con lati di lunghezza $a$, $b$, $c$, $d$ e diagonali di lunghezza $e$, $f$ è data dalla formula % lang-it

\begin{equation}
    S = \sqrt{(p-a)(p-b)(p-c)(p-d) - \frac 1 4 (ac + bd + pq)(ac + bd - pq)},
\end{equation}

gdzie $p$ oznacza raz jeszcze (i nie ostatni) połowę obwodu. % lang-pl

dove $p$ denota ancora una volta (e non per l'ultima) il semiperimetro. % lang-it

%